% !PRINTED 658
\weberpage{658}

Die in diesen Formeln vorkommenden Resolventen der Kreistheilung
$(r^{m_2}, \eta)$ sind Kreistheilungszahlen in einem h\"oheren
K\"orper. Ihre $m^{\text{ten}}$ Potenzen sind aber im K\"orper
$\Omega_m$ enthaltene Zahlen, die durch die Substitution $(r, r^n)$
in
\[
(r^{n m_2}, \eta)^m
\]
\"ubergehen. Ebenso sind die Verbindungen
\weq{13}{(r^{n m_2}, \eta)^{-1} (r^{m_2}, \eta)^n}
im K\"orper $\Omega_m$ enthalten [\S.\,174, (9), (11)]. Diese
Resolventen sind specielle F\"alle der allgemeinen Resolvente
$(r, x_0)$. Aus (12) erhalten wir, wenn wir mit $\varphi_n,
\varepsilon_n$ die conjugirten Functionale zu $\varphi$ und
$\varepsilon$ bezeichnen
\weq{14}{\omega_n = \varepsilon_n \varphi_n^m \prod^{p} (r^{m_2 n}, \eta)^m.}
Diese Formel lehrt uns die wichtigsten Eigenschaften der
Functionale $\varphi_n$ und $\varepsilon_n$ kennen, zun\"achst:

7. \gesperrt{Die Functionale $\varphi_n^m$ sind mit Zahlen des
K\"orpers $\Omega_m$ associirt, geh\"oren also der Hauptclasse an
(\S.\,152).}

Nach 1. ist $\omega_n^{-1} \omega_1^n = \alpha^m$ die % [sic] source omits "in" before $\Omega_m$
$m^{\text{te}}$ Potenz einer Zahl $\Omega_m$. Da auch die
Verbindungen (13) Zahlen in $\Omega_m$ sind, so k\"onnen wir
\[
\prod^{p} (r^{m_2 n}, \eta)^{-1} (r^{m_2}, \eta)^n = \frac{\alpha}{\beta}
\]
setzen, und erhalten aus (14)
\[
\varepsilon_n^{-1} \varepsilon_1^n (\varphi_n^{-1} \varphi_1^n)^m = \beta^m,
\]
worin $\beta$ eine Zahl in $\Omega_m$ ist.

Daraus geht hervor, dass das Product
\[
\beta \varphi_n \varphi_1^{-n} = \varepsilon
\]
ein Einheitsfunctional des K\"orpers $\Omega_m$ ist, und daraus
ergeben sich f\"ur die Functionale $\varphi_n$ und die
Einheitsfunctionale $\varepsilon_n$ die folgenden beiden S\"atze:

8. \gesperrt{Die conjugirten Functionale $\varphi_n$ haben die
Eigenschaft, dass alle Producte}
\[
\varphi_n \varphi_1^{-n}
\]
\gesperrt{der Hauptclasse angeh\"oren.}

9. \gesperrt{Die Einheitsfunctionale $\varepsilon_n$ haben die
Eigenschaft, dass die Producte}
\[
\varepsilon_n \varepsilon_1^{-n} = \varepsilon^m
\]
\gesperrt{$m^{\text{te}}$ Potenzen von Einheitsfunctionalen sind.}
